% SR-GNN：单个匿名会话的物品转移图（教学示例：登山包、水壶、登山包、相机）。
% 单栏图（figure），内容宽约78mm；回访合并为同一节点，出边按频次归一化。
\begin{tikzpicture}[x=1mm,y=1mm,
  font=\figurefont\footnotesize,
  gnode/.style={rectangle,draw=FigureModel!70,fill=FigureModel!10,align=center,
    text=BookInk,rounded corners=1.4pt,minimum height=9mm,
    text width=16mm,inner sep=2pt},
  glocal/.style={rectangle,draw=FigureOutput!70,thick,fill=FigureOutput!14,align=center,
    text=BookInk,rounded corners=1.4pt,minimum height=9mm,
    text width=16mm,inner sep=2pt},
  gflow/.style={-{Stealth[length=1.8mm,width=1.2mm]},draw=FigureOutput!85,line width=.85pt},
  gback/.style={-{Stealth[length=1.8mm,width=1.2mm]},draw=FigureModel!80,line width=.8pt,dashed},
  glab/.style={font=\figurefont\scriptsize,text=BookMuted,fill=white,inner sep=1pt}
]
  % 节点：回访的登山包只建一次
  \node[gnode] (a) at (16,20) {登山包\\（回访合并）};
  \node[gnode] (b) at (58,33) {水壶};
  \node[glocal] (c) at (58,9) {相机\\（末次局部）};

  % 转移边：包→水壶，水壶→包（回访），包→相机
  \draw[gflow] (a.north east) .. controls (36,34) .. (b.west);
  \node[glab,align=center] at (37,30) {出边 $1/2$};
  \draw[gback] (b.south west) .. controls (36,25) .. (a.north east);
  \node[glab,align=center] at (36,22) {回访};
  \draw[gflow] (a.south east) .. controls (38,12) .. (c.west);
  \node[glab,align=center] at (38,9) {出边 $1/2$};

  % 局部与全局说明
  \node[glab,align=left] at (16,-2)
    {局部：末次节点 $c$\\全局：$\alpha_k$ 对全部节点聚合};
\end{tikzpicture}
